\documentclass{article}

\usepackage{iclr2027_conference,times}
\usepackage{amsmath,amssymb,amsthm}
\usepackage{graphicx}
\usepackage{microtype}
\usepackage{hyperref}
\usepackage{url}

\newtheorem{definition}{Definition}
\newtheorem{lemma}{Lemma}
\newtheorem{theorem}{Theorem}
\newtheorem{proposition}{Proposition}
\newtheorem{corollary}{Corollary}

\newcommand{\preflow}{\textnormal{\textsc{PREFLOW}}}
\newcommand{\fcfs}{\textsc{FCFS}}
\newcommand{\edf}{\textsc{EDF}}
\newcommand{\slo}{\textsc{SLO}}
\newcommand{\isl}{\textsc{ISL}}
\newcommand{\ttft}{\textsc{TTFT}}
\newcommand{\R}{\mathbb{R}}

\title{\preflow: Safe Prefill Preemption\\
with FCFS-Relative Completion Guarantees}

% Authors are hidden while \iclrfinalcopy remains commented out.
\author{AUTHORS}

% Uncomment only for the camera-ready version.
\iclrfinalcopy

\begin{document}

\maketitle

\begin{abstract}
Production LLM traffic is both bursty and heterogeneous: prompts with orders of
magnitude different prefill costs can accumulate in the same queue.  Chunking
makes long prefills preemptible, but leaves a central policy question open:
\emph{when is it safe to let short work overtake them?}  FCFS avoids starvation
but causes head-of-line blocking; shortest-job policies improve mean latency but
provide no bound on the delay imposed on an individual request.  We present
\preflow, a scheduler that makes preemption safe and interpretable.  At arrival,
request $i$ receives an endogenous deadline
$D_i=a_i+\rho B_i$, where $B_i$ is its predicted completion latency under a
legal FCFS continuation from the observed state and $\rho\geq1$ is the maximum
allowed completion inflation.  A triangular attention-work model represents
both chunk progress and reusable cached prefixes.  At each chunk boundary,
\preflow{} serves the
highest-value request whose execution preserves joint feasibility of all fixed
deadlines.  On a preemptive single-resource model, this decision reduces to an
incremental EDF prefix-slack test that states exactly how much overtaking is
safe.  Under conservative work estimates, we prove
$(C_i-a_i)/B_i\leq\rho$ for every request, independent of future arrivals.  We
also extend the obligation across the KV-memory admission boundary using a
conservative drain envelope.  \preflow{} thus exploits preemption aggressively
without heuristic starvation controls or synthetic per-length SLOs, while
bounding only the delay introduced by scheduling rather than congestion already
present at arrival.
\end{abstract}

\section{Introduction}
\label{sec:introduction}

Production traces make queue order consequential.  BurstGPT documents rapidly
varying request concurrency in Azure OpenAI services, while Mooncake reports a
timestamped Kimi workload with a broad input-length distribution and substantial
prefix reuse~\citep{wang2024burstgpt,qin2024mooncake}.  The same prefill queue can
therefore mix short chat turns with long RAG, document, and agent contexts whose
attention costs differ by orders of magnitude.  In a prefill--decode
disaggregated deployment, a single long prompt can block an arriving burst at a
dedicated prefill worker.

Chunked prefill exposes preemption points and removes the need to execute a
prompt atomically~\citep{agrawal2024sarathi,agrawal2024medha,
hsieh2026flowprefill}.  Yet preemptibility alone is not a scheduling policy.
FCFS preserves arrival order but suffers head-of-line blocking.  Least-
remaining-work policies remove convoys, but a stream of short arrivals can
repeatedly preempt a long request.  Fixed bypass limits, aging multipliers, and
priority boosts can mitigate starvation, but their parameters do not express a
completion guarantee and interact unpredictably with load and request-size
mix.  The open problem is therefore not merely how to preempt, but how to
identify \emph{exactly when preemption remains safe}.

Existing SLO-aware schedulers answer a related but different question.  If a
request carries a meaningful external deadline---from a contract, an
interactive threshold, or a larger workflow---then the scheduler should use
it.  When no such deadline exists, however, synthesizing one from an offline
length profile creates a calibration problem.  The same 64K prompt may arrive
to an idle worker or behind seconds of queued work.  A length-bucket deadline
that is tight in the first state can be infeasible in the second; one loose
enough for the burst provides little protection when the system is idle.
Conditioning the table on load, mix, burstiness, hardware, and queue state
progressively learns an offline approximation to the state-dependent quantity
that is available directly at arrival.

\preflow{} derives this quantity from a concrete counterfactual.  Request $i$
receives the immutable deadline
\begin{equation}
    D_i = a_i + \rho B_i,
    \label{eq:deadline-intro}
\end{equation}
where $a_i$ is its arrival, $B_i$ is its predicted FCFS response from the
arrival state, and $\rho\geq1$ is the allowed inflation.  Thus $\rho=1.75$
permits at most $75\%$ more modeled response time than the arrival-time FCFS
baseline.  The deadline contracts when the worker is idle and expands when a
burst has already made low absolute latency impossible.  Unlike an arbitrary
SLO set, the constraints come with a feasible witness; unlike an urgency score,
they bound the additional harm caused by every future scheduling decision.

At each chunk boundary, \preflow{} ranks requests by a soft efficiency score,
then executes the first candidate that preserves all fixed deadlines.  This
separates \emph{what is profitable} from \emph{what is safe}: the score can be
changed without changing the protection contract.  In the single-resource
model, safety reduces to an EDF prefix-slack test.  It measures the exact amount
of service that a later-deadline request may consume before some protected
prefix becomes infeasible; when no such slack remains, the corresponding EDF
work must run.

This missing contract aligns with the interface exposed by current serving
stacks.  As Section~\ref{sec:related} documents, prominent routers choose a
prefill endpoint from cache affinity and load, then spill to another worker once
the preferred endpoint is saturated; SLO-aware modes, when present, commonly
translate fleet-level TTFT and TPOT targets into capacity.  The local prefill
queue therefore has request work, cache, and backlog state, but not necessarily
a meaningful deadline for every request.  \preflow{} turns exactly that state
into a contract: one global $\rho$ self-calibrates to each arrival, while the
feasibility shield provides an explicit bound without a workload-specific SLO
table.

This paper makes four contributions:
\begin{itemize}
    \item We introduce an endogenous, arrival-state completion contract that
    separates inherited congestion from scheduler-induced delay.
    \item We derive a preemptive scheduler that greedily optimizes a soft
    efficiency objective subject to exact EDF prefix feasibility, and show how
    much overtaking each decision can safely consume.
    \item We prove the per-request bound
    $(C_i-a_i)/B_i\leq\rho$, independent of future arrivals, under conservative
    work estimates.
    \item We extend the obligation across memory-constrained admission with a
    conservative KV-capacity drain envelope.
\end{itemize}

\section{Model and objective}
\label{sec:model}

\subsection{Cache-aware triangular work}

Request $i$ arrives at time $a_i$ with a prompt of $n_i$ tokens.  Let
$0\leq c_i<n_i$ be the reusable cached prefix observed at arrival and let
$h_i(t)$ be the longest prefix already available at time $t$, so
$h_i(a_i)=c_i$.  The prefill completes at $C_i$ when $h_i(C_i)=n_i$.

At an attention layer, each new token attends over its preceding history.
Ignoring a common constant, the cumulative work of a prefix is the triangular
number $W(n)$.  The cache-aware isolated and remaining work are therefore
\begin{equation}
    W(n)=\frac{n(n+1)}{2},\qquad p_i=W(n_i)-W(c_i),\qquad r_i(t)=W(n_i)-W(h_i(t)).
    \label{eq:work-model}
\end{equation}
An $x$-token chunk starting after history $h$ consumes
\begin{equation}
    \Delta W(h,x)=W(h+x)-W(h)=hx+\frac{x(x+1)}{2}.
    \label{eq:chunk-work}
\end{equation}
Thus the same simple model covers full prefills, partial progress, and prefix
reuse: a cache hit removes the triangular work already performed for the shared
prefix while retaining the attention from each new suffix token to that prefix.
Execution is non-preemptive within a chunk and preemptive at chunk boundaries.

Prefix caching is a target use case, not an exclusion.  If blocks credited at
arrival are evicted while $i$ waits, its work estimate increases.  A practical
implementation can recompute the affected queued work and correct its deadline
before the next decision, just as cache-aware JCT schedulers recalibrate waiting
requests~\citep{du2025prefillonly}.  The formal bound assumes that credited
blocks remain available or that their possible recomputation is charged
conservatively.  We treat dynamic deadline correction as a systems extension
and evaluate how often it is required.

\subsection{Compute and admission state}

An LLM serving engine cannot necessarily begin every arrived prefill.  We
partition the state into pending requests $P(t)$ and in-flight requests $I(t)$
whose execution state is resident.  An admission is legal only if its KV-cache
allocation respects memory capacity $M$.  vLLM exposes waiting and running
collections and admits against available KV blocks~\citep{kwon2023vllm,
vllmSchedulerDocs}; SGLang and TensorRT-LLM implement different forms of
running-batch and capacity management~\citep{zheng2024sglang,
tensorrtExecutorDocs}.  The names and reclamation mechanisms differ, but a
pending/in-flight boundary is common.

We first derive the scheduler on a unit-speed, preemptive work resource where
all admitted work is immediately executable.  Section~\ref{sec:admission}
extends the model to memory-constrained admission.

\subsection{Choosing among safe requests}

The hard constraint developed below tells us which chunks are safe, but often
leaves several choices.  We use intrinsic stretch to choose among them.  The
flow time of request $i$ is its arrival-to-completion latency
$F_i=C_i-a_i$; dividing by its cache-aware isolated work $p_i$ makes that
latency comparable across prompt sizes:
\begin{equation}
    S_i=\frac{F_i}{p_i}=\frac{C_i-a_i}{p_i}.
\end{equation}
Minimizing total stretch is weighted flow minimization with weight $1/p_i$.
With known remaining work, this motivates the weighted-shortest-remaining-job
priority
\begin{equation}
    \pi_i(t)=\frac{1}{p_i r_i(t)}.
    \label{eq:soft-score}
\end{equation}
The factor $1/p_i$ favors requests that are intrinsically small after cache
reuse, while $1/r_i(t)$ also recognizes a long request that is close to
completion.  This score expresses only a preference: another latency, utility,
or tenant objective can replace it without changing the safety test.  The hard
boundary uses $B_i$, not $p_i$, because protection must also reflect the backlog
that request $i$ inherited at arrival.

\section{FCFS-relative completion constraints}
\label{sec:constraints}

\subsection{Arrival-time baseline}

Freeze the system state at arrival $a_i$ and assume no future requests arrive.
Let $B_i$ be an upper bound on the time until $i$ completes under a legal,
work-conserving FCFS continuation from that state.

In the single-resource model,
\begin{equation}
    B_i=Q_i(a_i)+p_i,
    \label{eq:baseline-simple}
\end{equation}
where $Q_i(a_i)$ is all remaining work already present ahead of $i$.  In the
general model, $B_i$ is obtained from an FCFS reference simulation that also
accounts for memory release and admission.  This witness-based definition is
more robust than assuming that GPU batch times are perfectly additive.

We distinguish two dimensionless quantities:
\begin{equation}
    S_i=\frac{C_i-a_i}{p_i}
    \quad\text{and}\quad
    I_i=\frac{C_i-a_i}{B_i}.
    \label{eq:two-normalizers}
\end{equation}
$S_i$ is intrinsic stretch and defines the efficiency objective.  $I_i$ is
FCFS-relative completion inflation and defines protection.  The baseline $B_i$
is fixed at arrival; future requests cannot retroactively enlarge it.

\subsection{Constrained optimization}

For inflation limit $\rho\geq1$, assign $D_i$ by
Equation~\eqref{eq:deadline-intro}.  The conceptual offline problem is
\begin{align}
    \min_{\text{legal schedule}}\quad
        &\sum_i \frac{C_i-a_i}{p_i}
        \label{eq:optimization}\\
    \text{s.t.}\quad
        &C_i\leq D_i=a_i+\rho B_i && \forall i,\nonumber\\
        &\text{chunk, compute, and memory constraints.}&&\nonumber
\end{align}
\preflow{} is an online greedy solution: Equation~\eqref{eq:soft-score}
approximates the objective, while a feasibility invariant enforces the hard
constraints.

\subsection{Why the baseline is not an SLO table}

$B_i$ is conditioned on the resource state seen by request $i$, not only on its
length.  Thus two equal-length requests in the same trace may receive different
deadlines: one arrives to an idle worker, while the other inherits a burst.  A
per-length p99 table instead averages over the queue states, mix, hardware, and
execution policy used to construct it.  It can be a useful tuned baseline, but
does not retain the same meaning under workload or system shifts.

External SLOs remain appropriate when they encode application value.  They can
drive admission, provisioning, or SLO-goodput scheduling; if infeasible, the
system should reject work or record a violation.  For accepted requests without
such a deadline, \preflow{} provides the complementary scheduler-level
contract: it bounds added interference after accounting for inherited
congestion.

\section{PREFLOW scheduling}
\label{sec:scheduler}

\subsection{EDF prefix feasibility}

At a decision time $t$, order all unfinished, released requests by fixed
deadline:
\begin{equation}
    D_{(1)}\leq D_{(2)}\leq\cdots\leq D_{(m)},
\end{equation}
breaking ties by arrival order.  Define the residual slack of prefix $k$ as
\begin{equation}
    \sigma_k(t)=D_{(k)}-t-\sum_{j=1}^{k}r_{(j)}(t).
    \label{eq:prefix-slack}
\end{equation}
For released jobs on a preemptive unit-speed server, the state is feasible if
and only if $\sigma_k(t)\geq0$ for every $k$.  Necessity follows from the
capacity available before $D_{(k)}$; sufficiency follows from EDF.  We include
the standard argument in Appendix~\ref{app:prefix-proof}.

\begin{proposition}[Safe-chunk criterion]
\label{prop:safe-chunk}
Suppose request $i$ has EDF rank $m_i$.  Executing $q\leq r_i(t)$ units of $i$
preserves residual feasibility if and only if
\begin{equation}
    q\leq \min_{k<m_i}\sigma_k(t),
    \label{eq:safe-chunk}
\end{equation}
where the minimum of the empty set is $+\infty$.
\end{proposition}

\begin{proof}
For $k<m_i$, request $i$ is absent from prefix $k$: time advances by $q$ while
its remaining work does not decrease, hence $\sigma'_k=\sigma_k-q$.  For
$k\geq m_i$, the same prefix contains $i$: both available time and remaining
work decrease by $q$, hence $\sigma'_k=\sigma_k$.  Therefore exactly the
earlier prefixes can become infeasible, yielding
Equation~\eqref{eq:safe-chunk}.
\end{proof}

The proposition makes the scheduler incremental.  A preferred request spends
only the slack of requests with earlier deadlines.  The earliest-deadline
request is always safe in a feasible state.  If a proposed chunk is too large,
the engine can reduce its size to the available slack or defer it.

\subsection{Greedy feasible selection}

At every chunk boundary, \preflow{} applies the following rule:
\begin{enumerate}
    \item Rank legal request actions by decreasing soft score $\pi_i(t)$, with
    deterministic FCFS tie-breaking.
    \item For each candidate in that order, predict the work $q$ of its next
    chunk and test the successor state's deadline feasibility.
    \item Execute the first safe candidate.  If no preferred candidate is safe,
    execute an earliest-deadline action.
    \item Recompute scores, remaining work, and feasibility at the next chunk
    boundary.
\end{enumerate}

This policy preempts when doing so improves the soft objective and remains
feasible; otherwise it serves the tightest EDF \emph{prefix}, not a permanently
locked job.  Proposition~\ref{prop:safe-chunk} turns the protection rule into an
interpretable preemption budget: a later-deadline action may consume at most the
minimum slack of every prefix it bypasses.  A deadline-ordered tree can maintain
prefix sums and minima in logarithmic time; a linear scan suffices for modest
queues.

\section{Memory-constrained admission}
\label{sec:admission}

\subsection{The admission gap}

The compute-only result assumes that the selected request can run.  In an LLM
engine, a request can reach the head of $P(t)$ while insufficient KV capacity
prevents admission.  A guarantee that it merely reaches the first pending
position is therefore not a completion guarantee.

The exact generalization is state based.  A state contains remaining work,
$P(t)$, $I(t)$, resident allocations, capacity $M$, and all fixed deadlines.  It
is feasible if some legal sequence of admissions, chunks, and releases meets
every deadline.  \preflow{} may execute the highest-scoring action whose
successor remains feasible.  This definition yields a clean theorem but may be
expensive to evaluate, motivating a conservative sufficient test.

\subsection{Drain envelopes}

For pending request $i$, let $G_i(t)$ upper-bound the time needed to free enough
resident memory to admit $i$ once the scheduler commits to protecting it.  A
concrete drain policy may (i) stop admitting younger requests that increase the
blocking set, (ii) prioritize resident blockers whose completion releases the
required memory, and (iii) admit $i$ as soon as its reservation fits.  If
$r_i(t)$ upper-bounds its post-admission work, define the latest protection time
\begin{equation}
    L_i(t)=D_i-G_i(t)-r_i(t).
    \label{eq:latest-protection}
\end{equation}
By $L_i(t)$, the scheduler must make $i$ admission-critical and begin the drain
plan.  This formalizes the required headroom: $i$ must reach the protected
pending position early enough for the current in-flight set to drain and for
$i$ itself to execute.

The arrival baseline can equivalently be decomposed as
\begin{equation}
    B_i=Q_i^{\mathrm{pending}}(a_i)+G_i(a_i)+p_i.
    \label{eq:admission-baseline}
\end{equation}
Including $G_i(a_i)$ in $B_i$ is necessary, but the scheduler must also prevent
later admissions from enlarging $G_i$ after the deadline is fixed.

A conservative admission shield maintains four invariants: (1) once $i$ is
critical, do not admit younger work that increases $G_i$; (2) execute its
blocker-drain plan within the charged envelope; (3) admit $i$ as soon as its
reservation fits; and (4) after admission, preserve compute-prefix feasibility.
If several requests are critical, they are protected by deadline order.  Full
input-length reservation gives a simple, conservative realization; paging,
swapping, or destructive preemption can instead be represented as explicit
release actions with charged costs.

\section{Completion guarantee}
\label{sec:guarantee}

We state the result for a generic feasibility oracle.  The compute-only prefix
test is exact; the memory-aware shield is a sufficient implementation when its
drain estimates are conservative.

\paragraph{Assumptions.}
(A1) Every realized chunk duration is at most its charged work, and any cached
prefix credited at deadline assignment remains available or its possible
recomputation is included in that charge.
(A2) The scheduler reconsiders decisions at chunk boundaries and charges a full
non-preemptive chunk before dispatch.
(A3) At arrival, $B_i$ upper-bounds completion of $i$ in a legal FCFS witness
that does not violate older deadlines.
(A4) Every $G_i$ upper-bounds the protected drain policy, and later actions
cannot enlarge a protected blocking set.
(A5) \preflow{} executes only successor states accepted by the exact oracle or
a sound sufficient test.
(A6) The scheduler makes progress whenever unfinished work exists.

\begin{lemma}[Arrival feasibility]
\label{lem:arrival}
If the state immediately before $a_i$ is feasible, adding request $i$ with
$D_i=a_i+\rho B_i$ for $\rho\geq1$ preserves feasibility.
\end{lemma}

\begin{proof}
By A3, a legal continuation meets every old deadline and completes $i$ by
$a_i+B_i$.  Since $a_i+B_i\leq a_i+\rho B_i=D_i$, that continuation is a
feasible witness for the enlarged state.
\end{proof}

\begin{lemma}[Action invariance]
\label{lem:action}
If the current state is feasible and \preflow{} executes a safe action, the
successor state is feasible.
\end{lemma}

\begin{proof}
This follows from the definition of a safe action.  In the single-resource
model, Proposition~\ref{prop:safe-chunk} is an exact constructive test.  Under
memory constraints, A4--A5 make the conservative admission shield sound.
\end{proof}

\begin{theorem}[FCFS-relative completion bound]
\label{thm:bound}
Under A1--A6, every accepted request $i$ completes by its fixed deadline:
\begin{equation}
    C_i\leq D_i=a_i+\rho B_i,
    \qquad
    \frac{C_i-a_i}{B_i}\leq\rho.
    \label{eq:main-bound}
\end{equation}
\end{theorem}

\begin{proof}
The empty system is feasible.  Lemma~\ref{lem:arrival} preserves feasibility at
each arrival, and Lemma~\ref{lem:action} preserves it after each action.  A1 and
A4 ensure that realized compute and admission consume no more time than the
feasibility model reserves; A2 prevents an uncharged chunk overrun; and A6
excludes indefinite idling.  Thus the maintained feasible continuation can be
realized and no fixed deadline is missed.
\end{proof}

\begin{corollary}[Independence from future arrivals]
Once request $i$ is accepted, later arrivals cannot weaken
Equation~\eqref{eq:main-bound}.
\end{corollary}

The deadline and baseline are immutable.  Each later arrival must itself
preserve feasibility, and all subsequent actions remain subject to the existing
deadlines.  Future load may force \preflow{} to follow EDF more often, reducing
its freedom to optimize mean stretch, but it cannot legally consume $i$'s
reserved capacity.

\section{Discussion}
\label{sec:discussion}

\paragraph{Interpretation under bursts.}
The guarantee is intentionally relative.  A request arriving into a severe
burst may have a large $B_i$, because no scheduler can erase work that is
already present.  Nevertheless, the factor $\rho$ bounds the additional harm
caused by future overtaking.  This separates inherited congestion from
scheduler-induced slowdown.

\paragraph{Choosing $\rho$.}
A smaller $\rho$ approaches FCFS and leaves little room for short work to pass;
a larger $\rho$ improves stretch efficiency but permits more tail inflation.
The parameter has a direct operational meaning: $\rho=1.75$ permits at most
$75\%$ modeled completion inflation over the request's arrival-time FCFS
baseline.  A single global value yields the cleanest admission rule.
Class-specific values are possible if each arrival is rejected or relaxed
whenever its proposed deadline cannot be appended feasibly.

\paragraph{External SLOs.}
\preflow{} is not an argument against real application deadlines.  A
contractual tier or workflow deadline specifies value outside the scheduler and
should be honored whenever admission and provisioning make it feasible.  The
FCFS-relative deadline answers a different question for requests without such
deadlines: after accounting for inherited state, how much additional harm may
the scheduler create?  Benchmark-derived per-length p99 targets are useful
tuned baselines, but should not be conflated with external requirements.

\paragraph{Prediction error.}
Theorem~\ref{thm:bound} is exact in modeled work units.  A deterministic
wall-clock claim requires upper-confidence costs or a known multiplicative
error envelope.  If actual action time is at most $(1+\eta)$ times the point
estimate, charging $(1+\eta)$ times estimated work recovers the modeled bound.
Online correction can debit measured overruns from slack, but cannot repair one
unmodeled overrun larger than the remaining reserve.

\paragraph{Dynamic prefix-cache state.}
Prefix reuse changes work rather than the scheduling argument.  Stable or
pinned hits fall directly under Theorem~\ref{thm:bound}.  If an unexpected
eviction increases a queued request's work, the implementation can rebuild the
affected work and deadline before the next chunk boundary; the resulting
contract is relative to the corrected baseline.  Preserving the original
arrival-time contract instead requires pinning the credited blocks or reserving
their possible recomputation.  We expect the distinction to matter only for
requests whose observed hits are evicted while they wait, and report that
frequency separately.

\paragraph{Batching and non-additivity.}
GPU batch time is not perfectly additive.  The exact abstraction treats a
batch as a state transition with a predicted duration and a vector of request
progress.  A simpler conservative implementation can serialize accounting
while physically batching compatible work.  This may reject some safe
schedules, but does not invalidate the guarantee.

\paragraph{Scope.}
The theorem does not promise low absolute latency, competitiveness with an
offline optimum, or service for a request that cannot fit in an empty system.
It bounds prefill completion rather than full decode completion, and excludes
unbounded cost error, hardware failure, and permanent capacity loss.

\section{Related work}
\label{sec:related}

\paragraph{From preemption mechanisms to an ordering contract.}
Paged KV management and chunked prefill make it possible to reconsider a request
before its prompt completes~\citep{kwon2023vllm,agrawal2024sarathi,
zheng2024sglang}.  LoongServe adds elastic sequence parallelism, while Medha
pushes preemptive execution to million-token contexts through adaptive chunking,
Sequence Pipeline Parallelism, KV-Cache Parallelism, and multi-prefill
batching~\citep{wu2024loongserve,agrawal2024medha}.  These systems establish
\emph{how} long prefills can be split and executed efficiently.  \preflow{}
starts from that scheduling freedom and asks a separate question: which
reorderings may the worker perform without imposing unbounded harm?

\paragraph{Work-aware and SLO-aware scheduling.}
Known prefill work naturally motivates job-completion-time policies.
PrefillOnly serves single-token-output workloads with SRJF, continuously
recomputing each waiting request's JCT from its current prefix-cache hit and
subtracting a waiting-time offset for fairness~\citep{du2025prefillonly}.  It is
important precedent for treating prefix reuse as dynamic work, but its fairness
control remains a soft priority adjustment rather than a completion bound.

Medha is the closest comparison because its LARS scheduler explicitly combines
long-context preemption, heterogeneity, and deadlines~\citep{agrawal2024medha}.
LARS assumes that request $i$ arrives with a relative deadline $d_i$ and ranks
the normalized slack
$(a_i+d_i-t-w_i(t))/w_i^{\mathrm{total}}$.  This is a sensible policy when
$d_i$ encodes a valid per-request SLO.  It does not, however, derive that target
from the backlog encountered at arrival or test whether the selected chunk
preserves joint feasibility of the residual deadline set.  FlowPrefill likewise
uses finer preemption to maximize externally defined TTFT-SLO
goodput~\citep{hsieh2026flowprefill}; when targets conflict, goodput optimization
decides which requests miss them.  \preflow{} addresses the complementary case
in which no meaningful individual deadline is supplied: it constructs a
feasible target from the observed state and bounds the interference created by
all later reorderings.

\paragraph{Classical foundation.}
SRPT minimizes mean flow time when job sizes are known, but its treatment of
large jobs motivates the literature on slowdown and stretch
\citep{bender1998flow,muthukrishnan1999online,bansal2001srpt}.  EDF feasibility
for released preemptive jobs is also classical.  The novelty of \preflow{} is
the bridge between them: an online FCFS counterfactual creates each deadline,
and EDF prefix feasibility turns that deadline set into an exact safety
certificate for a modular efficiency policy.

\paragraph{Fleet routing and within-worker control.}
PD systems such as Splitwise, DistServe, and Mooncake use TTFT and TPOT
objectives to provision pools and assign work across prefill and decode
resources~\citep{patel2024splitwise,zhong2024distserve,qin2024mooncake}.  At
runtime, prominent open-source routers expose a cache/load interface rather than
requiring a deadline on every request.  vLLM Router provides sticky hashing and
power-of-two load balancing; llm-d prefers cache-resident endpoints, bypasses a
preferred endpoint after a calibrated prefill-throughput threshold marks it
saturated, and then selects by token load; Dynamo exposes active-prefill-token
busy thresholds.\footnote{Official descriptions:
\href{https://vllm.ai/blog/2025-12-13-vllm-router-release}{vLLM Router},
\href{https://github.com/llm-d/llm-d/blob/main/guides/precise-prefix-cache-routing/README.md}{llm-d precise prefix routing}, and
\href{https://docs.nvidia.com/dynamo/v-0-9-1/components/router/router-guide}{Dynamo Router}.}
Dynamo's optional SLA planner similarly translates TTFT/ITL targets into fleet
capacity, while its default throughput mode scales on queue depth and KV
utilization.\footnote{\href{https://docs.nvidia.com/dynamo/dev/knowledge-base/modular-components/planner/overview}{NVIDIA Dynamo Planner}.}
The common operational pattern is therefore cache-affine saturation followed by
spillover; per-request deadlines can be added, but are not a universal input to
the local prefill scheduler.  This is precisely where \preflow{} is attractive:
it needs only state already visible at the worker and one interpretable
inflation factor, yet gives every accepted request an explicit guarantee.

To our knowledge, prior LLM serving work has not used an online FCFS
counterfactual to construct per-request completion bounds and enforce them as a
joint feasibility invariant at every preemption point.  This is the core
novelty of \preflow{}; the work score, EDF condition, and chunking mechanism are
established components placed in service of that contract.

\section{Conclusion}
\label{sec:conclusion}

Bursts and service-time heterogeneity make preemption valuable; repeated
overtaking makes it dangerous.  \preflow{} resolves this tension by deriving a
fixed completion entitlement from the state each request encounters and using
joint feasibility to decide exactly when a preemption is safe.  It can optimize
any suitable soft priority while guaranteeing, under conservative compute and
admission models, that every request completes within $\rho$ times its
arrival-state FCFS baseline regardless of future arrivals.  This provides a
direct, load-adaptive alternative to heuristic starvation controls and
synthetic SLO calibration, at the within-prefill-worker scheduling layer
already exposed by modern PD serving stacks.

\bibliography{iclr2027_conference}
\bibliographystyle{iclr2027_conference}

\clearpage
\appendix

\section{Detailed prefix-feasibility proof}
\label{app:prefix-proof}

For released preemptive jobs on a unit-speed server, let $J_k$ be the first $k$
jobs in deadline order and
\begin{equation}
    R_k(t)=\sum_{j\in J_k}r_j(t).
\end{equation}
Any feasible schedule must process $R_k(t)$ units during $[t,D_{(k)}]$, whose
capacity is $D_{(k)}-t$.  Thus $R_k(t)\leq D_{(k)}-t$, equivalent to
$\sigma_k(t)\geq0$.

Conversely, run EDF.  If EDF first misses $D_{(k)}$, then over the relevant busy
interval it processed only jobs with deadlines no later than $D_{(k)}$;
otherwise this could not be the first missed deadline in a released,
work-conserving instance.  The required work of that prefix must exceed the
available interval, contradicting its nonnegative slack.  Hence the prefix
conditions are sufficient.

Now execute $q$ units of rank-$m$ request $i$.  For $k<m$, $i\notin J_k$, so
$R'_k=R_k$ while the interval loses $q$ units: $\sigma'_k=\sigma_k-q$.  For
$k\geq m$, $i\in J_k$, so both $R_k$ and the interval decrease by $q$:
$\sigma'_k=\sigma_k$.  This proves Proposition~\ref{prop:safe-chunk}.

\section{Examples}
\label{app:examples}

\paragraph{Long request followed by short work.}
At $t=0$, request $A$ arrives with $p_A=B_A=100$.  For $\rho=1.75$, its
deadline is $D_A=175$.  Suppose that at $t=1$, after one unit of progress,
$r_A=99$ and short jobs with higher soft scores begin arriving.  The singleton
prefix slack is
\begin{equation}
    \sigma_A(t)=175-t-99=76-t.
\end{equation}
Short work may consume at most 75 additional units after $t=1$.  Once the slack
is zero, every later-deadline action is unsafe and $A$ must run.  An unbounded
number of subsequent arrivals cannot change $D_A$.

\paragraph{Jointly tight prefix.}
At $t=5$, suppose $(D_A,r_A)=(20,8)$ and $(D_B,r_B)=(25,10)$.  Then
\begin{equation}
    \sigma_1=20-5-8=7,
    \qquad
    \sigma_2=25-5-(8+10)=2.
\end{equation}
A later-deadline request may consume only two work units because the joint
prefix $\{A,B\}$ is tighter than $A$ alone.  Pairwise bypass counters or
individual urgency scores do not directly expose this interaction.

\paragraph{Memory-blocked request.}
A pending request has $D_A=200$ and $r_A=40$.  Releasing enough resident KV
capacity takes at most $G_A=30$ under the protected drain plan.  Its latest
protection time is $L_A=200-30-40=130$.  Before 130, the scheduler may exploit
slack.  At 130, it must freeze harmful younger admissions, drain blockers,
admit $A$ by 160, and reserve its 40 units of service.  Moving $A$ to the
waiting-queue head at 175 would not suffice.

\section{Implementation-oriented feasibility interface}
\label{app:interface}

A serving engine can expose the theoretical shield through three operations:
\begin{enumerate}
    \item \textsc{PredictAction}$(s,u)$ returns a conservative duration and
    successor resource state for chunk or admission action $u$.
    \item \textsc{IsSafe}$(s')$ checks compute-prefix slack and all active drain
    obligations in successor state $s'$.
    \item \textsc{Fallback}$(s)$ returns an earliest-deadline compute action or
    the blocker-release action required by the earliest admission obligation.
\end{enumerate}
The priority layer is then engine agnostic: it sorts candidates by
Equation~\eqref{eq:soft-score} and selects the first action accepted by
\textsc{IsSafe}.  Engine-specific logic is confined to state prediction,
memory release, and legal chunk construction.

\end{document}
